\chapter{Homotopy limits and colimits: the theory}

Recall that since $\varprojlim $ and $\varinjlim $ are not necessarily homotopic $\mathcal{M} ^{\mathcal{D} } \to \mathcal{M} $, we want to compute their derived functors. Since colimits are left adjoints, we try to compute $\mathbb{L} \varinjlim $; respectively, $\mathbb{R} \varprojlim $. The choice of derived functor does not matter if $\mathcal{M} $ is saturated (the localization inverts precisely those maps which are weak equivalences), which it always is in practice (for example, in all model categories).

We will see later in the text that our current method is not the only answer; there are approaches via finding a Reedy categorical structure on $\mathcal{D} $, or by taking $\mathcal{M} $ to be a cofibrantly generated or combinatorial model category. The former are fairly specialized, and the latter are computationally difficult, so we specialize instead in simplicial model categories $\mathcal{M} $ by producing deformations from the bar and cobar constructions. We use the properties axiomatized by \cref{lma:3.7} to do this. Next chapter explores examples of interesting simplicial model categories, and later we explore ``local'' universality onditions of the homotopy (co)limit.

\section{The homotopy limit and colimit functors}

We write $\hocolim : \mathcal{M} ^{\mathcal{D} }  \to \mathcal{M} $ for the left derived functor of $\varinjlim $, and $\holim $ for the right derived functor of $\varprojlim $. We start by fixing at least one construction satisfying this property on the point-set level. This chapter will be dedicated to showing the following theorems, which define a specific construction of homotopy (co)limits. We first prove the relevant corollary, since the theorem is rather long.

\begin{theorem}\label{thm:5.1}
	Let $\mathcal{M} $ be a simplicial model category with cofibrant replacement $Q$ and fibrant replacement $R$\footnote{Every arrow admits a factorization as a weak equivalence followed by a fibration; in particular, if this fibration is to the terminal object, we find that the object is weakly equivalent to a fibrant object. Since model categories in particular have this factorization be functorial, we obtain a fibrant replacement functor $R : \mathcal{M}  \to \mathcal{M} $ replacing objects with their fibrant replacement. Dually, we obtain a cofibrant replacement. Source: Hovey, $n$Lab}. The pair
	\[
		B(\mathcal{D} ,\mathcal{D} ,Q(-)): \mathcal{M} ^{\mathcal{D} }  \to \mathcal{M} ^{\mathcal{D} } \qquad
		\begin{tikzcd}
			B(\mathcal{D} ,\mathcal{D} ,Q(-))\ar[" \varepsilon _Q ", Rightarrow, r] & Q\ar[" q ", Rightarrow, r] & 1
		\end{tikzcd}
	\]
	is a left deformation for $\varinjlim : \mathcal{M} ^{\mathcal{D} }  \to \mathcal{M} $, where we recall our notation gives the functor
	\[
		F\mapsto B(\mathcal{D} ,\mathcal{D} ,Q(F))
	.\]
	Dually,
	\[
		C(\mathcal{D} ,\mathcal{D} ,R(-)): \mathcal{M} ^{\mathcal{D} }  \to \mathcal{M} ^{\mathcal{D} } \qquad
		\begin{tikzcd}
			1\ar[" r ", Rightarrow, r] & R\ar[" \eta_R ", Rightarrow, r] & C(\mathcal{D} ,\mathcal{D} ,R(-))
		\end{tikzcd}
	\]
	is a right deformation for $\varprojlim : \mathcal{M} ^{\mathcal{D} }  \to \mathcal{M} $.
\end{theorem}
\begin{corollary}\label{cor:5.1}
	If $\mathcal{M} $ is a simplicial model category and $\mathcal{D} $ is any small category, then the functors $\varinjlim ,\varprojlim : \mathcal{M} ^{\mathcal{D} }  \to \mathcal{M} $ admit left and derived functors, called $\hocolim $ and $\holim $, respectively, defined by
	\[
		\hocolim_\mathcal{D}  \coloneqq \mathbb{L} \varinjlim_\mathcal{D} \cong B(*,\mathcal{D} ,Q(-))\qquad \text{and} \qquad \holim_\mathcal{D} \coloneqq \mathbb{R} \varprojlim _\mathcal{D} \cong C(*,\mathcal{D} ,R(-))
	.\]
\end{corollary}
\begin{proof}
	Note here that $*$ denotes the constant functor at the terminal object in $\sSet$, so recall that
	\[
		*\otimes _\mathcal{D} - \coloneqq \int^{d\in \mathcal{D} } *\otimes -\cong \int^{d\in \mathcal{D} } -\cong \varinjlim -
	\]
	since $\Delta^0$ is also the unit for the monoidal (cartesian) product in $\sSet$. By \cref{thm:5.1}, the left derived functor of colim exists, and by \cref{thm:2.1} it is given as precomposing by the left deformation
	\[
		\mathbb{L} \varinjlim _\mathcal{D} (-)\cong \varinjlim_\mathcal{D} (-)\circ B(\mathcal{D} ,\mathcal{D} ,Q(-))=\varinjlim _\mathcal{D} B(\mathcal{D} ,\mathcal{D} ,Q(-))\cong *\otimes _\mathcal{D} B(\mathcal{D} ,\mathcal{D} ,Q(-))
	.\]
	Fubini's theorem for ends\footnote{this is imo the best name for a theorem in basic category theory, shoutout my man mac lane} gives us that functor tensor products commute past either variable in the two-sided bar construction. More precisely,
	\begin{align*}
		*\otimes _\mathcal{D} B(\mathcal{D} ,\mathcal{D} ,Q(-))&=\int^{d\in \mathcal{D} } *\otimes B(\mathcal{D} (-,d),\mathcal{D} ,Q(-))\\
									   &=\int^{d\in \mathcal{D} } *\otimes \int^{[n]\in \boldsymbol{\Delta}} \Delta^n \otimes B_n(\mathcal{D} ,(-,d),\mathcal{D} ,Q(-))\\
									   &\cong\int^{d\in \mathcal{D} } \int^{[n]\in \boldsymbol{\Delta} } *\otimes \Delta^n \otimes B_n(\mathcal{D} (-,d),\mathcal{D} ,Q(-))\\
									   &\cong \int^{[n]\in \boldsymbol{\Delta} } \int^{d\in \mathcal{D} } *\otimes \Delta^n \otimes B_n(\mathcal{D} (-,d),\mathcal{D} ,Q(-))\\
									   &\cong \int^{[n]\in \boldsymbol{\Delta} } \int^{d\in \mathcal{D} } (*\times \Delta^n)\otimes B_n(\mathcal{D} (-,d),\mathcal{D} ,Q(-))\\
									   &\cong \int^{[n]\in\boldsymbol{\Delta} } \int^{d\in \mathcal{D} } (\Delta^n\times *)\otimes B_n(\mathcal{D} (-,d),\mathcal{D} ,Q(-))\\
									   &\cong \int^{[n]\in \boldsymbol{\Delta} } \int^{d\in \mathcal{D} } \Delta^n \otimes *\otimes \coprod_{\vec{d} : [m] \to \mathcal{D} } \mathcal{D} (d_n,d)\otimes Q(-)(d_0)\\
									   &\cong \int^{[n]\in\boldsymbol{\Delta} } \int^{d\in \mathcal{D} } \Delta^n \otimes \coprod_{\vec{d} : [m] \to \mathcal{D} } *\otimes \mathcal{D} (d_n,d)\otimes Q(-)(d_0)\\
									   &\cong\int^{[n]\in\boldsymbol{\Delta} } \int^{d\in \mathcal{D} } \Delta^n\otimes B(*\otimes \mathcal{D} (-,d),\mathcal{D} ,Q(-))\\
									   &\cong \int^{[n]\in\boldsymbol{\Delta} } \Delta^n \otimes B_n(*\otimes _\mathcal{D} \mathcal{D},\mathcal{D} ,Q(-))\\
									   &\cong B(*\otimes_\mathcal{D} \mathcal{D} ,\mathcal{D} ,Q(-))
	.\end{align*}
	The steps involve the fact that left adjoint functors preserve colimits, \cref{lma:3.5}, and definitions \ref{dfn:3.9} and \ref{dfn:3.11}. One could state that this suffices as well as a proof of \cref{lma:3.6}. All one needs to understand the proof is to make precise the domains and codomains of each functor, then simply work through the tensors, viewing sets as enriched over simplicial sets as discussed previously whenever needed. Hence
	\[
		*\otimes _\mathcal{D} B(\mathcal{D} ,\mathcal{D} ,Q(-))\cong B(*\otimes _\mathcal{D} \mathcal{D} ,\mathcal{D} ,Q(-))\cong B(*,\mathcal{D} ,Q(-))
	.\]
	This last isomorphism follows by the coYoneda lemma
	\[
		F(c)\cong \int^{d\in \mathcal{D} } \mathcal{D} (d,c)\cdot F(d)
	.\]
	We can view this as tensoring down into $\Set $, and being isomorphic to the cartesian product. Hence we have $*\otimes _\mathcal{D} \mathcal{D} $, evaluated at $d$, gives the constant functor $d\mapsto *$ by evaluation.

	The dual statement follows by the dual string of isomorphisms:
	\[
		\mathbb{R} \varprojlim _\mathcal{D} (-)\cong \varprojlim _\mathcal{D} C(\mathcal{D} ,\mathcal{D} ,R(-)\cong \left\{ *,C(\mathcal{D} ,\mathcal{D} ,R(-) \right\} ^{\mathcal{D} } 
	.\]
	The cotensor is an end, so by the same logic it commutes through. However since cobar is contravariant in its first variable, limits become colimits, or more precisely the cotensor becomes a tensor:
	\[
		\left\{ *,C(\mathcal{D} ,\mathcal{D} ,R(-)) \right\} ^{\mathcal{D} } \cong C(*\otimes _\mathcal{D} \mathcal{D} ,\mathcal{D} ,R(-))\cong C(*,\mathcal{D} ,R)
	.\]
\end{proof}

With this fact established, \cref{thm:5.1} remains to be shown. We begin with a series of lemmas.

\begin{lemma}\label{lma:5.1}
	The natural weak equivalence $\varepsilon : B(\mathcal{D} ,\mathcal{D} ,F) \Rightarrow F$ makes $B(\mathcal{D} ,\mathcal{D} ,-)$ a left deformation for $\mathcal{M} ^{\mathcal{D} } $.
\end{lemma}
\begin{proof}
	Recall from the last part of chapter 4 that $B_{\bullet} (\mathcal{D} (-,d),\mathcal{D} ,F)$ admits an augmentation and extra degeneracy, which induces a natural map $\varepsilon : B(\mathcal{D} ,\mathcal{D} ,F) \Rightarrow F$ which is a pointwise simplicial homotopy equivalence. This is a nontrivial fact and I will gloss over it for now. and Since $\mathcal{M} $ is a simplicial model category, \cref{lma:3.7} makes this a pointwise weak equivalence. Hence the two-sided bar construction $B(\mathcal{D} ,\mathcal{D} ,-)$ is a deformation for $\mathcal{M} ^{\mathcal{D} } $, since we recall that is simply an endofunctor and a natural weak equivalence to 1.
\end{proof}

We then obtain the deformation in \cref{thm:5.1} by composing $\varepsilon $ with the left deformation $q : Q \Rightarrow 1$ to the cofibrant objects of the simplicial model category, which exists by \cref{lma:3.7}. This is certaintly a deformation, so it suffices to show it's indeed a deformation for the colimit functor. This means we must show that colimit preserves pointwise weak equivalences between diagrams in a full subcategory of $\mathcal{M} ^\mathcal{D} $ containing the image of $B(\mathcal{D} ,\mathcal{D} ,Q(-))$.

\begin{lemma}\label{lma:5.2}
	Let $(Q,q)$ be a left deformation on $\mathcal{M} $ and $F : \mathcal{M}  \to \mathcal{N} $. If
	\begin{itemize}
		\item $FQ$ is homotopical;
		\item $FqQ : FQ^2 \Rightarrow FQ$ is a natural weak equivalence;
	\end{itemize}
	then $(Q,q)$ is a left deformation for $F$, meaning $F$ is homotopical on a full subcategory of cofibrant objects if $F$ preserves weak equivalences in the image of $Q$ and also preserves the weak equivalences $qQ$.
\end{lemma}
\begin{proof}
	Let $x,y\in\mathcal{M}$. The obvious thing to do is to show weak equivalences in $\mathcal{M} (Q(x),Q(y))$ are preserved by $F$. For any $f : Q(x) \to Q(y)$, we have that $FQ(f)$ is a weak equivalence, since $FQ$ is homotopical. Naturality of $FqQ$ gives
	\[\begin{tikzcd}[row sep = 2.5em, column sep = 2.5em]
		FQ^2(x)\ar[" FQ(w) ", r] \ar[" F(q_{Q(x)}) "', d]  & FQ^2(y)\ar[" F(q_{Q(y)} ) ", d]  \\
	FQ(x)\ar[" F(w) "', r]  & FQ(y)
	\end{tikzcd}\]
	This reads as the equality
	\[
		F(w)\circ F(q_{Q(x)} )=F(q_{Q(y)} )\circ FQ(w)
	.\]
	Since $FqQ$ is a natural weak equivalence, we have that $F(w)$ is a weak equivalence by the 2-of-3 property: the right side is a weak equivalence since weak equivalences are closed under composition, and $F(q_{Q(x)} )$ is a weak equivalence.
\end{proof}

By \cref{lma:5.2}, it suffices to show that $B(*,\mathcal{D} ,Q(-))\cong \varinjlim _\mathcal{D} B(\mathcal{D} ,\mathcal{D} ,Q(-))$ is homotopical, and that the composite
\[\begin{tikzcd}[row sep = 2.5em, column sep = 2.5em]
\varinjlim _\mathcal{D} B(\mathcal{D} ,\mathcal{D} ,Q(B(\mathcal{D} ,\mathcal{D} ,Q(F))))\ar[" \varinjlim _\mathcal{D} \cdot \varepsilon  ", r]  & \varinjlim _\mathcal{D} Q(B(\mathcal{D} ,\mathcal{D} ,Q(F)))\ar[" \varinjlim _\mathcal{D} \cdot q ", r]  & \varinjlim _\mathcal{D} B(\mathcal{D} ,\mathcal{D} ,Q(F))
\end{tikzcd}\]
is a natural weak equivalence. We show this in the next section.

\section{Homotopical aspects of the bar construction}

To finish the proof, we must explore the homotopical properties of the bar construction, which follow from the fact that the two-sided simplicial bar is \emph{Reedy cofibrant} when $G$ and $F$ are pointwise cofibrant. A Reedy cofibrant simplicial object is one in which the inclusion of the degenerate $n$-simplicies into all $n$-simplicies is a cofibrant. A degenerate $n$-simplex is supposed to be a simplex in the image of a degeneracy map, but I have no idea what this means in non-concrete categories. Anyways the inclusion is called the latching map and the degenerate $n$-simplicies are called the $n$th latching object.

\begin{lemma}\label{lma:5.3}
		Let $\mathcal{D} $ be small and $\mathcal{M} $ be a simplicial model category. If $F : \mathcal{D}  \to \mathcal{M} $ is pointwise cofibrant, then $B_{\bullet} (*,\mathcal{D} ,F)$ is Reedy cofibrant.
\end{lemma}
\begin{proof}
	Recall $B_n(*,\mathcal{D} ,F)$ is the coproduct over $N(\mathcal{D} )_n$ of the image under $F$ of the first object in the composable sequence of arrows. Hence, the $n$th latching object sits inside $B_n(*,\mathcal{D} ,F)$ as the coproduct indexed by degenerate simplicies in $N(\mathcal{D} )_n$. By a lemma proved in 6 chapters, the cofibrations, and hence the cofibrant objects, are closed under coproducts. By the same lemma, cofibrations are closed under pushout, and hence coproduct inclusions of cofibrant objects are also cofibrations. By hypothesis, every object in one of these coproducts is cofibrant, so the above observations implythat the $n$th latching map is a cofibration.
\end{proof}

I have no idea what this means I copied it directly from the book. I might review it when I learn what these words mean. For now, note this explains the role of the functor $Q$ in the bar construction. Postcomposition by $Q$ yields a pointwise cofibrant replacement of the diagram $F$, so the theorem will apply.

Recall that tensoring with any simplicial set preserves cofibrant objects. Thus, $B_{\bullet} (G,\mathcal{D} ,F)$ is Reedy cofibrant by the same argument for any pointwise cofibrant $F$ with codomain in a simplicial model category, and $G : \mathcal{D} ^{\mathrm{op}}  \to \sSet$. We now introduce a standard result that will this time be proven in 9 chapters, which we need to prove the current result.

\begin{theorem}\label{thm:5.2}
	If $\mathcal{M} $ is a simplicial model category, then
	\[
		\left| - \right| : \mathcal{M} ^{\boldsymbol{\Delta} ^{\mathrm{op}} }  \to \mathcal{M} 
	\]
	is left Quillen with respect to the Reedy model structure. In particular, $\left| - \right| $ sends Reedy cofibrant simplicial objects to cofibrant objects and preserves pointwise weak equivalences between them.
\end{theorem}

\begin{remark}
	This is the strongest result we can give. It is not true that geometric realization preserves all pointwise homotopy equivalences, only those between cofibrant objects. This is why it is not a homotopy colimit always, since it is not in general homotopical.
\end{remark}

\begin{corollary}\label{cor:5.2}
	The functor $B(G,\mathcal{D} ,-)$ preserves weak equivalences between pointwise cofibrant objects. In particular, $B(*,\mathcal{D} ,Q(-))$ is homotopical. Dually, when $F$ is pointwise cofibrant, $B(-,\mathcal{D} ,F)$ preserves weak equivalences.
\end{corollary}
\begin{proof}
	Since $B_{\bullet} (G,\mathcal{D},F)$ is Reedy cofibrant for $F$ pointwise cofibrant, by \cref{thm:5.2} so is $B(G,\mathcal{D} ,F)$. In particular, pointwise weak equivalences between these objects are preserved, so $B(G,\mathcal{D} ,-)$ preserves the relevant weak equivalences between pointwise cofibrant objects. Since $Q$ is a left deformation into cofibrant objects, we obtain that $B(*,\mathcal{D} ,Q(-))$ is homotopical.
\end{proof}

This shows the first objective of the proof. Now we must show the relevant natural map is a weak equivalence. Naturality gives
\[\begin{tikzcd}[row sep = 2.5em, column sep = 2.5em]
B(\mathcal{D} ,\mathcal{D} ,Q(B(\mathcal{D} ,\mathcal{D} ,Q(F))))\ar[" \varepsilon_{QB}  ", r] \ar[" B(\mathcal{D} \text{,} \mathcal{D} \text{,} q) "', d]  & Q(B(\mathcal{D} ,\mathcal{D} ,Q(F)))\ar[" q ", d]  \\
B(\mathcal{D} ,\mathcal{D} ,B(\mathcal{D} ,\mathcal{D} ,Q(F)))\ar[" \varepsilon _B "', r]  & B(\mathcal{D} ,\mathcal{D} ,Q(F))
\end{tikzcd}\]
By the discussion following \cref{lma:5.3}, $B_{\bullet} (\mathcal{D} ,\mathcal{D} ,Q(F))$ is pointwise Reedy cofibrant, since $Q$ is a deformation into cofibrant objects. By \cref{thm:5.2}, $B(\mathcal{D} ,\mathcal{D} ,Q(F))$ is pointwise cofibrant. Taking the colimit of the left side of the diagram gives
\[
	B(*,\mathcal{D},q): B(*,\mathcal{D} ,Q(B(\mathcal{D} ,\mathcal{D} ,Q(F)))) \to B(*,\mathcal{D} ,B(\mathcal{D} ,\mathcal{D} ,F))
.\]
\cref{cor:5.2} says that $B(*,\mathcal{D} ,-)$ preserves weak equivalences between pointwise cofibrant diagrams, hence this map is a weak equivalences. It remains only to be shown that the image of the bottom map under colim is a weak equivalence, by the 2-of-3 property.
\begin{lemma}\label{lma:5.4}
	When $F$ is pointwise cofibrant, the functor $\varinjlim _\mathcal{D} $ preserves the weak equivalence
	\[
		\varepsilon_{B(\mathcal{D} ,\mathcal{D} ,F)} : B(\mathcal{D} ,\mathcal{D} ,B(\mathcal{D} ,\mathcal{D} ,F)) \to B(\mathcal{D} ,\mathcal{D} ,F)
	.\]	
\end{lemma}
\begin{proof}
	Recall that $\varepsilon_F : B(\mathcal{D} ,\mathcal{D} ,F) \to F$ is the geometric realization of a map between simplicial objects determined by the maps
	\[
		\mathcal{D} (d_n,d)\times \left( \prod_{i=0}^{n-1} \mathcal{D} (d_i,d_{i+1} )\right)\cdot F(d_0)\to F(d)
	.\]
	These maps compose the chains of arrows in $\mathcal{D} $ to obtain a map $d_0\to d$, then evaluate this at $F$ to obtain $F(d_0)\to F(d)$. Since functor tensor products commute with both variables of the two-sided bar construction, and since $D(-,d)\otimes _\mathcal{D} F\cong F(d)$, we have
	\[
	B(\mathcal{D} ,\mathcal{D} ,F)\cong B(\mathcal{D} ,\mathcal{D} ,\mathcal{D} ,\otimes _\mathcal{D} F)\cong B(\mathcal{D} ,\mathcal{D} ,\mathcal{D} )\otimes _\mathcal{D} F\xlongrightarrow{\varepsilon \otimes _\mathcal{D} 1} \mathcal{D} \otimes _\mathcal{D} F\cong F
	.\]
	Here, $\varepsilon $ is the geometric realization of the map of simplicial objects
	\[
		\varepsilon : B_{\bullet} (\mathcal{D} (-,d),\mathcal{D} ,\mathcal{D} (d',-)) \to \mathcal{D} ,(d',d)
	\]
	defined by composing the arrows. This augmented simplicial object admits both forward and backwards contracting homotopies, which shows that $G\otimes _\mathcal{D} \varepsilon $ is a natural weak equivalence. The argument is then computed by a commutative diagram that I don't want to copy down.
\end{proof}

